\documentclass[11pt]{article}
\usepackage[a4paper,margin=18mm]{geometry}
\usepackage[no-math]{fontspec}
\setmainfont{Latin Modern Roman}
\usepackage{amsmath,amssymb,amsthm,xfrac,xspace,hyperref,graphicx,comment}
\excludecomment{DefTralics}
\providecommand{\dpl}{\displaystyle}
\providecommand{\backmatter}{}
\providecommand{\loccit}{\emph{loc. cit.}}

\providecommand{\bibliofont}{\small}
\newcommand{\ie}{i.e.,\xspace}
\newcommand{\eg}{e.g.,\xspace}
\newcommand{\etc}{etc.\xspace}
\newcommand{\cf}{cf.\xspace}
\newcommand{\resp}{resp.\xspace}
\newcommand{\smaller}{\small}
\newcommand{\Subsection}{\subsection}
\newcommand{\Subsubsection}{\subsubsection}
\newcommand{\skpt}{\smallskip}
\emergencystretch=3em
\setcounter{section}{1}
% Begin literal retained input author-preamble.tex
\usepackage{mathrsfs}
\let\mathcal\mathscr

% Literal retained local macro bbm
\RequirePackage{fontspec}
\newfontfamily\jepblackboardfont[ItalicFont={Latin Modern Math},BoldFont={Latin Modern Math},BoldItalicFont={Latin Modern Math}]{Latin Modern Math}
\makeatletter
\expandafter\def\csname jepbb@1\endcsname{\text{\jepblackboardfont\upshape\char"1D7D9}}
\DeclareRobustCommand{\mathbbm}[1]{\csname jepbb@#1\endcsname}
\makeatother



\graphicspath{{istrati-et-al_fig/}}
\newcommand{\Changel}{\mathcode`l="7160}
\newcommand{\Changelback}{\mathcode`l="716C}
\newcommand{\NoChangel}[1]{%
\expandafter\let\csname old\string#1\endcsname=#1
\let#1=\relax
\newcommand{#1}{\mathcode`l="716C\csname old\string#1\endcsname\mathcode`l="7160 }%
}

\NoChangel{\log}
\NoChangel{\ln}
\NoChangel{\lim}
\NoChangel{\limsup}
\NoChangel{\liminf}
\NoChangel{\varlimsup}
\NoChangel{\varliminf}
\NoChangel{\varinjlim}
\NoChangel{\varprojlim}

\newenvironment{enumeratei}
{\bgroup\def\theenumi{\roman{enumi}}\def\theenumii{\arabic{enumii}}\begin{enumerate}}
{\end{enumerate}\egroup}

\newenvironment{smallpmatrix}{\big(\begin{smallmatrix}}{\end{smallmatrix}\big)}

\newtheorem{theorem}{Theorem}[section]
\newtheorem{proposition}[theorem]{Proposition}
\newtheorem{lemma}[theorem]{Lemma}
\newtheorem{corollary}[theorem]{Corollary}

\theoremstyle{definition}
\newtheorem{definition}[theorem]{Definition}
\newtheorem{remark}[theorem]{Remark}
\newtheorem{example}[theorem]{Example}

\newcommand{\Psfrac}[2]{(\sfrac{#1}{#2})}
\newcommand{\psfrac}[2]{\sfrac{(#1)}{#2}}
\newcommand{\spfrac}[2]{\sfrac{#1}{(#2)}}
\newcommand{\pspfrac}[2]{\sfrac{(#1)}{(#2)}}

\newcommand{\sbullet}{{\raisebox{1pt}{$\scriptscriptstyle\bullet$}}}

\let\wt\widetilde
\let\sharp\#
\let\epsilon\varepsilon

\DeclareMathOperator{\codim}{codim}
\DeclareMathOperator{\St}{St}

\newcommand\from{\mathchoice{\longleftarrow}{\leftarrow}{\leftarrow}{\leftarrow}}
\newcommand\mto{\mathchoice{\longmapsto}{\mapsto}{\mapsto}{\mapsto}}
\newcommand\hto{\mathrel{\lhook\joinrel\to}}
\newcommand{\To}[1]{\mathchoice{\xrightarrow{\textstyle\kern4pt#1\kern3pt}}{\stackrel{#1}{\longrightarrow}}{}{}}


\let\al\alpha
\let\be\beta
\let\la\lambda
\let\del\partial
\let\ov\overline
\let\sminus\setminus
\let\un\underline

\newcommand{\BB}{\mathbb{B}}
\newcommand{\CC}{\mathbb{C}}
\newcommand{\DD}{\mathbb{D}}
\newcommand{\NN}{\mathbb{N}}
\newcommand{\PP}{\mathbb{P}}
\newcommand{\QQ}{\mathbb{Q}}
\newcommand{\RR}{\mathbb{R}}
\newcommand{\Ss}{\mathbb{S}}
\newcommand{\TT}{\mathbb{T}}
\newcommand{\ZZ}{\mathbb{Z}}

\newcommand{\Oo}{\mathcal{O}}

\newcommand{\e}{\mathrm{e}}

\newcommand{\ui}{\mathfrak{i}}

\newcommand{\one}{\mathbbm{1}}

\DeclareMathOperator{\Aut}{Aut}
\DeclareMathOperator{\Bl}{Bl}
\DeclareMathOperator{\coker}{coker}
\DeclareMathOperator{\GL}{GL}
\DeclareMathOperator{\Hom}{Hom}
\DeclareMathOperator{\id}{\mathsf{id}}
\DeclareMathOperator{\im}{im}
\DeclareMathOperator{\Imm}{Im}\let\Im\Imm
\DeclareMathOperator{\Int}{Int}
\DeclareMathOperator{\Inv}{Inv}
\DeclareMathOperator{\orb}{orb}
\DeclareMathOperator{\ord}{ord}
\DeclareMathOperator{\Pic}{Pic}
\DeclareMathOperator{\Ree}{Re}\let\Re\Ree
\DeclareMathOperator{\Spec}{Spec}

\newcommand{\abs}[1]{\vert #1\vert}
\newcommand{\annfan}[1]{\hat{#1}^0}
\newcommand{\transp}[1]{{#1}^{\scriptscriptstyle\mathrm{T}}}
\newcommand{\ce}{\mathcal{C}^\infty}
\newcommand{\Ka}{Kähler\xspace}

\usepackage{xparse}   %to use commands with multiple options

\NewDocumentCommand{\colfan}{ +m +m +m +O{} }
{
\Sigma_{#1}^{#2 \,{}^\centerdot #3}
}
\NewDocumentCommand{\colvar}{ +m +m +m +O{} }
{
\wt{X}_{#1}^{#2 \,{}^\centerdot #3}
}
\NewDocumentCommand{\colmap}{ +m +m +m +O{p} }
{
{#4}_{#1}^{#2 \centerdot #3}
}
\NewDocumentCommand{\idemap}{ +m +m +O{} }
{
\sigma{#3}_{#1}^{#2}
}

% End literal retained input author-preamble.tex

\begin{document}
\begin{center}{\Large Stable item 2108}\end{center}
\paragraph{Source and attribution.}
Nicolina Istrati, Alexandra Otiman, Massimiliano Pontecorvo, Matteo Ruggiero, \emph{Toric Kato manifolds}, Journal de l’École polytechnique — Mathématiques 9 (2022), publisher version, DOI 10.5802/jep.208. \href{https://jep.centre-mersenne.org/articles/10.5802/jep.208/}{Publisher article}. Authors' text: \href{https://creativecommons.org/licenses/by/4.0/}{CC BY 4.0}.
\paragraph{Editorial problem boundary.}
Prove only implication(2) to(3) of Theorem10.3 below for the exact Kato data and corresponding compact Kato manifold. A Kahler universal cover forces the printed compactified modified ball to be projective. Retain the complete compatible embeddings of iterated fundamental domains, punctured-ball construction and regularized-maximum gluing lemma; Miyaoka extension and Kahler-Moishezon projectivity remain the original permitted references. Other equivalences, toric classification, metric nonexistence and SVG illustrations are outside the selected claim.
\paragraph{Difficulty (editorial): H2.}
Miyaoka extension and Moishezon projectivity do not transfer the universal-cover Kahler form to the modified ball on their own. Proposition2.4 constructs compatible holomorphic embeddings by iterating the Kato fundamental domains and proving that they exhaust the punctured modified ball. Lemma10.2 then glues the pulled-back form to a Fubini–Study multiple through a controlled regularized maximum, providing the new construction needed before the final projectivity theorem applies.
\paragraph{Editorial transcription note.}
Original author language, mathematics and ordering are preserved. Publisher front matter is replaced by this independent article wrapper. Bibliography is recovered from matching cached publisher HTML, including its TeX formulas. No new proof or mathematical repair is supplied. Surrounding original results are retained for conservative context and are not extra selected corpus claims.
\clearpage

% Begin literal retained input author-body.tex
\section{Preliminaries on Kato manifolds}\label{prel}

In the present paper, a Kato manifold will be a compact complex manifold obtained as follows \cite{kato}, \cite{dl84}. We let $\BB:=\{(z_1, \ldots, z_n) \in \CC^n | |z_1|^2+\ldots+|z_n|^2 < 1\}$ be the standard open ball in $\CC^n$. Consider a modification $\pi: \hat{\BB} \to \BB$ at $0$ and a holomorphic embedding $\sigma: \overline{\BB} \hookrightarrow \hat{\BB}$. Glue the two boundaries of $\hat{\BB} \sminus \sigma(\BB)$ via the real analytic CR-diffeomorphism $\gamma:=\sigma \circ \pi$. The result is a compact complex manifold $X(\pi, \sigma)$, named here a \textit{Kato manifold}. Any small neighborhood of the image of $\del\hat\BB$ in $X(\pi,\sigma)$ is a global spherical shell. We shall refer to the couple $(\pi, \sigma)$ as \emph{Kato data} and to $F:=\pi \circ \sigma : (\BB, 0) \to (\BB, 0)$ as \emph{the corresponding germ}.

The class of Kato manifolds introduced in \cite{kato} is slightly more general, as $\pi$ is allowed to be a modification at more than one point, so~that the resulting manifolds are proper modifications of the ones we described above \cite[Lem.\,2.7, Part\,I]{dl84}. However, for the present discussion this generality will not make any difference.

We denote by $q:\wt{X}\to X$ the universal cover of a Kato manifold $X=X(\pi,\sigma)$. The deck group $\Gamma$ of $q$ is canonically isomorphic to $\ZZ$, and we will denote also by $\gamma$ the positive generator of $\Gamma$, which is indeed induced by the map $\sigma\circ\pi$.

We start by settling several general facts that will be needed in the paper.

\begin{lemma}\label{interFB}
Let $F:(\BB,0)\to (\BB,0)$ be a holomorphic germ with $F(\ov\BB)\subset\BB$. Then $\bigcap_{m\in\NN}F^m(\BB)=\{0\}$.
\end{lemma}
\begin{proof}
By hypothesis there exists $0<r<1$ such that $F(\BB)\subset\BB_r$. By the Schwarz lemma \cite[Th.\,6]{shabat}, we have $\|F(z)\|\leq r\|z\|$ for any $z\in\BB$, hence $\|F^m(z)\|\leq r^m$ for any $m>0$ and any $z\in\BB$. This implies that $\lim_{m\to\infty}\sup_{z\in\BB}\|F^m(z)\|=0$, and hence $\mathrm{diam}(\bigcap_{m\in\NN}F^m(\BB))=0$. The conclusion then follows.
\end{proof}

Given a map $F:(\CC^n,0)\to (\CC^n,0)$, define its stable set:
\begin{equation*}\textstyle
W^s(F)=\{z\in\CC^n\mid \lim_{m\to\infty}\|F^m(z)\|=0\}.
\end{equation*}

\begin{lemma}\label{reuniuneFB}
For any holomorphic map $F:(\CC^n,0)\to (\CC^n,0)$ with $F(\ov \BB)\subset\BB$, we~have $W^s(F)=\bigcup_{m\in\ZZ}F^m(\BB)$, where for $m=-k<0$, $F^{-k}(\BB)$ is the preimage of~$\BB$ via $F^k$.
\end{lemma}
\begin{proof}
The previous lemma shows that $\BB\subset W^s(\BB)$. Clearly $F^m(\BB)\subset W^s(F)$ for any $m>0$ and also for any $m<0$. Conversely, if $z\in W^s(F)$ then there exists $m>0$ such that $F^m(z)\in\BB$, which shows the desired equality.
\end{proof}

Suppose that $F:(\CC^n,0)\to(\CC^n,0)$ is given by a Kato data $(\pi,\sigma)$, $F=\pi\circ\sigma$, such that $\pi$ is a modification at $0$.
For any $k>0$, let $H_k:=F^{-k}(0)$ so~that $H_k\subseteq H_{k+1}$ and put $H_\infty:=\bigcup_{k>0} H_k$. Let us set $\Inv(F):=\CC^n\sminus H_\infty$ and $W^*(F):=W^s(F)\cap\Inv(F)$.

\begin{proposition}\label{stableF}
Let $F:(\CC^n,0)\to(\CC^n,0)$ be a holomorphic map with $F(0)=0$ given by a Kato data $(\pi,\sigma)$ via $F=\pi\circ\sigma$, so~that $\pi$ is a modification at $0$ with exceptional divisor $E$. Suppose that $H_\infty=H_m$ for some $m>0$. Let $D$ be the divisor on $X(\pi,\sigma)$ induced by $E$ after gluing. Then we have a biholomorphism $X(\pi,\sigma)\sminus D\cong W^*(F)/\langle F\rangle$.
\end{proposition}
\begin{proof}
If we put $\BB^*:=\BB\cap\Inv(F)$, then by \ref{reuniuneFB} we have $W^*(F)=\bigcup_{m\in\ZZ}F^m(\BB^*)$. Furthermore, by \ref{interFB}, we have:
\begin{equation*}
\bigcup_{m\in\ZZ}\left(F^m(\BB^*)\sminus F^{m+1}(\BB^*)\right)=\bigcup_{m\in\ZZ}F^m(\BB^*)\sminus \bigcap_{m\in\ZZ}F^m(\BB^*)=W^*(F).
\end{equation*}
Now it is clear that $D\cap \bigl(\hat\BB\sminus \sigma(\BB) \bigr)=\pi^{-1}(H_\infty)\cap (\hat\BB\sminus \sigma(\BB))$ and that
\[
\pi:\bigl(\hat\BB\sminus \sigma(\BB)\bigr)\sminus \pi^{-1}(H_\infty)\to \BB^*\sminus F(\BB^*)
\]
establishes a biholomorphism satisfying $\pi\circ\gamma=F\circ \pi$, where $\gamma=\sigma\circ\pi$. This implies in particular that the group $\langle F\rangle$ acts freely and properly on $W^*(F)$ and that $\pi$ induces the desired isomorphism $X(\pi,\sigma)\sminus D\cong W^*(F)/\langle F\rangle$.
\end{proof}

\begin{proposition}\label{embBall}
Let $(\pi:\hat\BB\to\BB,\,\sigma:\BB\to\hat\BB)$ be a Kato data with germ $F=\pi\circ\sigma$ and let $X=X(\pi,\sigma)$ be the corresponding Kato manifold. Then there exists a holomorphic open embedding $\phi:\hat\BB\sminus \{\sigma(0)\}\to \wt{X}$.
\end{proposition}

\Changel
\begin{proof}
We recall that $\wt{X}=\bigsqcup_{m\in\ZZ} W_m \big/\sim$, where $W_m=\hat\BB\sminus \sigma(\BB)$ for any $m\in\ZZ$ and $\sim$ indicates that $W_m$ is glued to $W_{m+1}$ via $\gamma=\sigma\circ\pi$.

Let us put, for any $l\geq 1$, $\wt{X}_l:=\bigsqcup_{1\leq m\leq l} W_m \big/\sim \, \,\subset \wt{X}$. Also, following \cite[\S 1]{iop}, denote by $(\pi_l:\hat\BB^{(l)}\to\BB,\sigma_l:\BB\to\hat\BB^{(l)})$ the Kato data obtained by composing $(\pi,\sigma)$ with itself $l$ times. Then it is easy to see that $\wt{X}_l=\hat\BB^{(l)}\sminus \sigma_l(\BB)$.

If we denote by $\pi_{(l)}:\hat\BB^{(l)}\to\hat\BB^{(l-1)}$ the induced map, then using that $\pi_{(l)}\circ\sigma_l=\sigma_{l-1}\circ F$ we find that $\pi^{-1}_{(l)}$ gives an embedding of $\hat\BB^{(l-1)}\sminus \sigma_{l-1}(F(\BB))$ into $\hat\BB^{(l)}\sminus \sigma_l(\BB)$. Therefore we find inductively, for any $l\geq 1$, open embeddings
\begin{equation*}
\phi_l:Q_l:=\hat\BB\sminus \sigma(F^{l-1}(\BB))\to\hat\BB^{(l)}\sminus \sigma_l(\BB)\subset \wt{X}
\end{equation*}
which satisfy $\phi_{l+1}|_{Q_l}=\phi_l$. Now since $\bigcap_{l\geq 1} F^l(\BB)=\{0\}$ by Lemma~\ref{interFB}, this means that $\bigcup_{l\geq 1} Q_l=\hat\BB\sminus \{\sigma(0)\}$, and thus the family $\{\phi_l\}_{l\geq 1}$ naturally defines an open embedding $\phi:\hat\BB\sminus \{\sigma(0)\}\to \wt{X}$, which concludes the proof.
\end{proof}


% End literal retained input author-body.tex

\setcounter{section}{9}

% Begin literal retained input author-core-proof.tex
\section{Hermitian geometry of Kato manifolds}\label{metrici}

Kato manifolds cannot admit Kähler metrics since their fundamental group is isomorphic to $\mathbb{Z}$. However, a large class of them carry \emph{locally conformally Kähler metrics (lcK)}, as shown in \cite[Th.\,2.2]{iop}. In this section, we give a characterization for the existence of lcK metrics on Kato manifolds. Furthermore, we investigate the possibility of endowing a Kato manifold with other special Hermitian metrics, such as balanced, pluriclosed, Hermitian symplectic or strongly Gauduchon.

Recall the following definition:
\begin{definition}\label{deflck} A \emph{locally conformally Kähler metric (lcK)} on a complex manifold~$X$ is a Hermitian metric $\Omega$ satisfying $d\Omega=\theta\wedge\Omega$ for a closed one-form $\theta$. If $\theta$ is not exact, $\Omega$ is called strictly lcK.
\end{definition}

This definition is equivalent to any of the following formulations:
\begin{enumerate}
\item There exists a covering with open sets $(U_i)_i$ of $X$ and smooth functions $f_i$ on every $U_i$ such that $e^{-f_i}\Omega$ is Kähler.
\item The universal cover $\wt{X}$ admits a Kähler metric $\wt{\Omega}$ such that $\gamma^*\wt{\Omega}=c_{\gamma}\wt{\Omega}$, $c_{\gamma}>0$, for any deck-transformation $\gamma$.
\end{enumerate}

For our characterization purpose, we start by giving a general principle used to modify Kähler metrics on modifications of balls, already used in \cite{bru} and \cite{iop}, which will be needed.

\begin{lemma}\label{extKA}
Let $\pi:\hat\CC^n\to \CC^n$ be a proper modification at $0$, let $\rho(z)=\rho(\|z\|)\in\ce(\CC^n,\RR)$ be a strictly plurisubharmonic function depending only on $\|z\|$, and let $\omega$ be a \Ka metric on $\hat{\ov\BB}:=\pi^{-1}(\ov\BB)$. Then for any $0<s<1$, there exists $\la>0$ and a \Ka metric $\wt \omega$ on $\hat\CC^n$ so~that
\begin{equation*}
\wt{\omega}|_{\hat\BB_s}=\omega, \quad\wt{\omega}|_{\hat\CC^n\sminus \hat\BB}=\la\cdot\pi^*dd^c\rho.
\end{equation*}
\end{lemma}
\begin{proof}
There exists a pluriharmonic function $\psi$ on $\BB$, smooth and strictly plurisubharmonic on $\BB\sminus \{0\}$, so~that $\pi_*\omega=dd^c\psi$. By the maximum principle, the functions
\begin{equation*}
u,v:(0,1]\to \RR, \quad u(t)=\max_{\del\BB_t}\rho, \quad v(t)=\max_{\del\BB_t}\psi
\end{equation*}
are strictly increasing. It follows that for any $0<s<1$,
\begin{gather*}
r:=u(s)\equiv \rho|_{\del\BB_s}<R:=u(1)\equiv \rho|_{\del\BB}, \quad m:=\min_{\del\BB_s}\psi\leq v(s)<M:=v(1).
\end{gather*}
Thus, for any $\al>1$, we can take
\begin{equation*}
\la:=\al\,\frac{M-m}{R-r}>0, \quad c\in\left(M-\la R,m-\la r\right)\neq\emptyset,
\end{equation*}
so~that $\la\rho+c>\psi$ on $\del\BB$ and $\la\rho+c<\psi$ on $\del\BB_s$. Therefore, if $\wt{\psi}$ is the function on~$\CC^n$ defined as the regularized maximum between $\psi$ and $\la\rho+c$ on $\BB\sminus \BB_s$, equal to $\psi$ on $\BB_s$ and equal to $\la\rho+c$ on $\CC^n\sminus \BB$, then the metric $\pi^*(dd^c\wt{\psi})$ on $\hat\CC^n\sminus \pi^{-1}(0)$ glues to $\omega$ to define the desired metric $\wt{\omega}$ on $\hat\CC^n$.
\end{proof}

Now we are ready to characterize the existence of lcK metrics. For this, let us note that any proper modification $\pi:\hat\BB\to\BB$ induces naturally proper modifications $\pi:\hat\CC^n\to\CC^n$ and $\pi:\hat{\CC\PP}^n\to\CC\PP^n$, where $\hat{\CC\PP}^n$ is the compactification of $\hat\CC^n$ with a hyperplane at infinity.

\begin{theorem}\label{lcK}
Let $(\pi:\hat\BB\to\BB,\sigma:\overline{\BB}\to\hat\BB)$ be a Kato data and let $X=X(\pi,\sigma)$ be the corresponding Kato manifold.
The following are equivalent:
\begin{enumerate}
\item\label{lcK1} $X$ admits an lcK metric;
\item\label{lcK2} $\wt{X}$ admits a \Ka metric;
\item\label{lcK3} $\hat{\CC\PP}^n$ is a projective manifold;
\item\label{lcK4} $\hat{\mathbb{C}}^n$ admits a Kähler metric.

\end{enumerate}
\end{theorem}

\begin{proof}
The implications $\eqref{lcK1}\Rightarrow\eqref{lcK2}$ and $\eqref{lcK3}\Rightarrow\eqref{lcK4}$ are clear, while the implication $\eqref{lcK4}\Rightarrow \eqref{lcK1}$ is precisely Brunella's theorem \cite{bru} adapted to all complex dimensions \cite[Th.\,2.2]{iop}. Indeed, in \cite[Th.\,2.2]{iop} we showed that if $\pi$ is a composition of smooth blow-ups, then $X$ is lcK. The same proof works if one merely supposes that~$\hat\BB$, or equivalently $\hat\CC^n$, is Kähler. We are thus left with showing $\eqref{lcK2} \Rightarrow \eqref{lcK3}$.

If $\wt{X}$ is Kähler, then Proposition~\ref{embBall} implies that $\hat\BB\sminus \{\sigma(0)\}$ is also \Ka, so~by Miyaoka's extension theorem \cite[Prop.\,A]{miyaoka}, $\hat\BB$ admits a Kähler metric $\omega$. Now putting $\rho:=\log(1+\|z\|^2)$ in Lemma~\ref{extKA}, we find that $\hat\CC^n$ admits a \Ka metric $\wt{\omega}$ which glues to $\la\cdot\pi^*\omega_{FS}$ on $\hat{\CC\PP^n}\sminus \hat\BB$ for some $\la>0$, where $\omega_{FS}$ is the standard Fubini-Study metric on $\CC\PP^n$. Thus we have shown that $\hat{\CC\PP^n}$ is \Ka. Since on the other hand $\hat{\CC\PP^n}$ is a modification of $\CC\PP^n$, it is Moishezon, and therefore $\hat{\CC\PP^n}$ is projective by Moishezon's theorem \cite[Chap.\,I, Th.\,11]{moi}. This concludes the proof.
\end{proof}

% End literal retained input author-core-proof.tex

\begin{thebibliography}{99}
\bibitem{berkovich} V. G. Berkovich - “A non-Archimedean interpretation of the weight zero subspaces of limit mixed Hodge structures” , in Algebra, arithmetic, and geometry: in honor of Yu. I. Manin. Vol. I , Progress in Math. , vol. 269 , Birkhäuser Boston, Boston, MA, 2009, p. 49-67
\bibitem{bis} J.-M. Bismut - “A local index theorem for non-Kähler manifolds” , Math. Ann. 284 (1989) no. 4, p. 681-699
\bibitem{boucksom-jonsson} S. Boucksom \& M. Jonsson - “Tropical and non-Archimedean limits of degenerating families of volume forms” , J. Éc. polytech. Math. 4 (2017), p. 87-139
\bibitem{bo} L. Battisti \& K. Oeljeklaus - “A generalization of Sankaran and LVMB manifolds” , Michigan Math. J. 64 (2015) no. 1, p. 203-222
\bibitem{b} F. Bosio - “Variétés complexes compactes: une généralisation de la construction de Meersseman et López de Medrano-Verjovsky” , Ann. Inst. Fourier (Grenoble) 51 (2001) no. 5, p. 1259-1297
\bibitem{bru} M. Brunella - “Locally conformally Kähler metrics on Kato surfaces” , Nagoya Math. J. 202 (2011), p. 77-81
\bibitem{bs} C. Bănică \& O. Stănăşilă - Algebraic methods in the global theory of complex spaces , Editura Academiei \& John Wiley \& Sons, Bucharest \& London-New York-Sydney, 1976
\bibitem{c} G. R. Cavalcanti - “Hodge theory of SKT manifolds” , Adv. Math. 374 (2020), article ID 107270, 42 pages
\bibitem{da} V. I. Danilov - “The geometry of toric varieties” , Uspekhi Mat. Nauk 33 (1978) no. 2, p. 85-134
\bibitem{dl84} G. Dloussky - Structure des surfaces de Kato , Mém. Soc. Math. France (N.S.) , Société Mathématique de France, Paris, 1984 no. 14
\bibitem{dot} G. Dloussky, K. Oeljeklaus \& M. Toma - “Class $\rm VII_0$ surfaces with $b_2$ curves” , Tôhoku Math. J. (2) 55 (2003) no. 2, p. 283-309
\bibitem{favre} C. Favre - “Degeneration of endomorphisms of the complex projective space in the hybrid space” , J. Inst. Math. Jussieu 19 (2020) no. 4, p. 1141-1183
\bibitem{fp} A. Fujiki \& M. Pontecorvo - “Anti-self-dual bihermitian structures on Inoue surfaces” , J. Differential Geometry 85 (2010) no. 1, p. 15-71
\bibitem{ft} A. Fino \& A. Tomassini - “Blow-ups and resolutions of strong Kähler with torsion metrics” , Adv. Math. 221 (2009) no. 3, p. 914-935
\bibitem{fulton} W. Fulton - Introduction to toric varieties , Annals of Math. Studies , vol. 131 , Princeton University Press, Princeton, NJ, 1993
\bibitem{gr} A. Grothendieck - “Sur quelques points d’algèbre homologique” , Tôhoku Math. J. (2) 9 (1957), p. 119-221
\bibitem{ega} A. Grothendieck - “Éléments de géométrie algébrique. IV. Étude locale des schémas et des morphismes de schémas. II” , Publ. Math. Inst. Hautes Études Sci. (1965) no. 24, p. 1-231
\bibitem{in1} M. Inoue - “New surfaces with no meromorphic functions” , in Proc. ICM (Vancouver, B.C., 1974), Vol. 1 , 1975, p. 423-426
\bibitem{in2} M. Inoue - “New surfaces with no meromorphic functions. II” , in Complex analysis and algebraic geometry (W. L. Baily \& T. Shioda, eds.), Iwanami Shoten \& Cambridge Univ. Press, Tokyo \& Cambridge, 1977, p. 91-106
\bibitem{iop} N. Istrati, A. Otiman \& M. Pontecorvo - “On a class of Kato manifolds” , Internat. Math. Res. Notices (2021) no. 7, p. 5366-5412
\bibitem{is} N. Istrati - “A characterisation of toric locally conformally Kähler manifolds” , J. Symplectic Geom. 17 (2019) no. 5, p. 1297-1316
\bibitem{kato} M. Kato - “Compact complex manifolds containing “global” spherical shells. I” , in Proc. of the Intern. Symposium on Algebraic Geometry (Kyoto Univ., Kyoto, 1977) , Kinokuniya Book Store, Tokyo, 1978, p. 45-84
\bibitem{kato79} M. Kato - “Some remarks on subvarieties of Hopf manifolds” , Tokyo J. Math. 2 (1979) no. 1, p. 47-61
\bibitem{kns} K. Kodaira, L. Nirenberg \& D. C. Spencer - “On the existence of deformations of complex analytic structures” , Ann. of Math. (2) 68 (1958), p. 450-459
\bibitem{mall91} D. Mall - “The cohomology of line bundles on Hopf manifolds” , Osaka J. Math. 28 (1991) no. 4, p. 999-1015
\bibitem{mic} M. L. Michelsohn - “On the existence of special metrics in complex geometry” , Acta Math. 149 (1982) no. 3-4, p. 261-295
\bibitem{miyaoka} Y. Miyaoka - “Extension theorems for Kähler metrics” , Proc. Japan Acad. 50 (1974), p. 407-410
\bibitem{moi} B. G. Moĭshezon - “On $n$-dimensional compact complex varieties with $n$ algebraically independent meromorphic functions. I” , in Seven papers on algebra, algebraic geometry and algebraic topology , Transl., Ser. 2 , vol. 63 , American Mathematical Society, Providence, RI, 1967, p. 51-93
\bibitem{naka83} I. Nakamura - “${\rm VII}_{0}$ surfaces and a duality of cusp singularities” , in Classification of algebraic and analytic manifolds (Katata, 1982) , Progress in Math. , vol. 39 , Birkhäuser Boston, Boston, MA, 1983, p. 333-378
\bibitem{naka84} I. Nakamura - “On surfaces of class ${\rm VII}_0$ with curves” , Invent. Math. 78 (1984) no. 3, p. 393-443
\bibitem{oda} T. Oda - Lectures on torus embeddings and applications. (Based on joint work with Katsuya Miyake.) , Lect. Math. Phys., Math., Tata Inst. Fundam. Res. , vol. 58 , Springer \& Tata Inst. of Fundamental Research, Berlin \& Bombay, 1978
\bibitem{oda2} T. Oda - Convex bodies and algebraic geometry , Ergeb. Math. Grenzgeb. (3) , vol. 15 , Springer-Verlag, Berlin, 1988
\bibitem{P9} D. Popovici - “Deformation limits of projective manifolds: Hodge numbers and strongly Gauduchon metrics” , Invent. Math. 194 (2013) no. 3, p. 515-534
\bibitem{P10b} D. Popovici - “Stability of strongly Gauduchon manifolds under modifications” , J. Geom. Anal. 23 (2013) no. 2, p. 653-659
\bibitem{P14} D. Popovici - “Deformation openness and closedness of various classes of compact complex manifolds; examples” , Ann. Scuola Norm. Sup. Pisa Cl. Sci. (5) 13 (2014) no. 2, p. 255-305
\bibitem{san} G. K. Sankaran - “A class of non-Kähler complex manifolds” , Tôhoku Math. J. (2) 41 (1989) no. 1, p. 43-64
\bibitem{shabat} B. V. Shabat - Introduction to complex analysis. Part II , Transl. of Math. Monogr. , vol. 110 , American Mathematical Society, Providence, RI, 1992
\bibitem{st} J. Streets \& G. Tian - “A parabolic flow of pluriclosed metrics” , Internat. Math. Res. Notices (2010) no. 16, p. 3101-3133
\bibitem{teissier} B. Teissier - “Two points of the boundary of toric geometry” , in Singularities, algebraic geometry, commutative algebra, and related topics , Springer, Cham, 2018, p. 107-117
\bibitem{tel} A. Teleman - “Donaldson theory on non-Kählerian surfaces and class VII surfaces with $b_2=1$” , Invent. Math. 162 (2005) no. 3, p. 493-521
\bibitem{tsu} H. Tsuchihashi - “Certain compact complex manifolds with infinite cyclic fundamental groups” , Tôhoku Math. J. (2) 39 (1987) no. 4, p. 519-532
\bibitem{yzz} S.-T. Yau, Q. Zhao \& F. Zheng - “On Strominger Kähler-like manifolds with degenerate torsion” , 2019
\end{thebibliography}
\end{document}
